\metatitle{Algebras and rings}
\metadescription{Brief definitions of Hecke, Temperley--Lieb, and Clifford algebras, with an overview of double affine Hecke algebras and their role in symmetric functions.}



\section[algebras]{Algebras appearing in representation theory and combinatorics}

\subsection[heckeAlgebra]{The Hecke algebra}

Let $(W,S)$ be a \hyperref[coxeterGroups]{Coxeter system}, with length
function $\ell$ and identity element $e$.  The \defin{Iwahori--Hecke algebra}
$\mathcal{H}_W$ is a deformation of the group algebra $\setZ[W]$.  One
common normalization defines $\mathcal{H}_W$ over $\setZ[v,v^{-1}]$ with
generators $H_s$, $s\in S$, satisfying the braid relations of $W$ and the
quadratic relation
\[
  (H_s-v)(H_s+v^{-1})=0.
\]
For each $w\in W$, the products along reduced words define a standard basis
$\{H_w : w\in W\}$.  At $v=1$, this specializes back to the group algebra.

\name{Jonathan Gruber} introduces pseudo-centralizer subalgebras of affine
Hecke algebras \cite{Gruber2026x}.  The construction starts inside the double
affine Hecke algebra and gives a subalgebra $\mathcal{B}$ of the affine Hecke
algebra which can be viewed as a $v$-deformation of Lam's affine
Fomin--Stanley subalgebra, a combinatorial model for the
\hyperref[schurAffine]{affine Grassmannian Schur-basis story}.  In type $A$
and in types $B_2$ and $G_2$, Gruber constructs a standard basis indexed by
minimal finite-Weyl-group coset representatives in the affine Weyl group, and
then a canonical basis defined using the Kazhdan--Lusztig basis.  The paper
also formulates positivity conjectures for this canonical basis and relates
the construction to the center of the affine Hecke algebra.
\name{Haruto Murata} proves a quantum imaginary Schur--Weyl duality for
quiver Hecke algebras of untwisted affine type $A$ \cite{Murata2026x}.
For arbitrary quiver-Hecke parameters, the resulting parameter $t$ determines
an Iwahori--Hecke algebra of the symmetric group, and the endomorphism
algebras in the imaginary strata are identified with these Hecke algebras.
The same framework computes characters of simple modules using dual canonical
bases and Kazhdan--Lusztig polynomials, and identifies standard-module
characters with PBW vectors under semisimplicity hypotheses.

\name{Naihuan Jing} and \name{Ning Liu} establish quantum super Schur--Weyl
duality for mirabolic Hecke algebras with $q$ indeterminate
\cite{JingLiu2026Mirabolic}. They derive a super Frobenius formula in
\hyperref[hookSchur]{hook Schur functions}, a Murnaghan--Nakayama rule,
and a Roichman character formula using super mirabolic RSK.

\subsection[kazhdanLusztigPolynomials]{Kazhdan--Lusztig polynomials}

\name{David Kazhdan} and \name{George Lusztig} introduced
\defin{Kazhdan--Lusztig polynomials} in their study of Coxeter groups and
Hecke algebras \cite{KazhdanLusztig1979}.  They are polynomials
$P_{u,w}(q)$ indexed by pairs $u,w\in W$, with $P_{u,w}(q)=0$ unless
$u\leq w$ in \hyperref[bruhatOrder]{Bruhat order} and $P_{w,w}(q)=1$.

The polynomials arise as transition coefficients between the standard basis of
$\mathcal{H}_W$ and the self-dual Kazhdan--Lusztig basis.  They encode deep
geometric and representation-theoretic information; for Weyl groups they are
closely related to intersection cohomology of Schubert varieties and to
characters in category $\mathcal{O}$.

Several variants occur elsewhere on this site, including parabolic
Kazhdan--Lusztig polynomials in the Schur expansion of
\hyperref[LLT]{LLT polynomials}, and matroid Kazhdan--Lusztig polynomials in
\hyperref[matroidKazhdanLusztigPolynomial]{matroid theory}.

\subsection[temperleyLiebAlgebra]{The Temperley--Lieb algebra}

Let $R$ be a commutative ring, $\delta\in R$, and $n\geq1$.
The \defin{Temperley--Lieb algebra} $\operatorname{TL}_n(\delta)$ is the
unital associative $R$-algebra generated by $U_1,\dotsc,U_{n-1}$ with relations
\begin{align*}
 U_i^2&=\delta U_i,\\
 U_iU_jU_i&=U_i &&\text{if }|i-j|=1,\\
 U_iU_j&=U_jU_i &&\text{if }|i-j|>1.
\end{align*}
Its diagram basis consists of noncrossing pairings of $n$ points on the top
and $n$ points on the bottom of a rectangle. We multiply diagrams by stacking
them and replacing each closed loop by a factor of $\delta$.
Thus the algebra is free of rank $C_n=\frac{1}{n+1}\binom{2n}{n}$ over $R$.
For an introduction to both presentations, see
\href{https://www.cs.ox.ac.uk/files/286/final.pdf}{S.~Abramsky's exposition, Section~2}.

With $\delta=v+v^{-1}$, the assignment $H_{s_i}\mapsto v-U_i$ gives a
surjection from the type $A$ \hyperref[heckeAlgebra]{Hecke algebra} onto
$\operatorname{TL}_n(\delta)$, in the normalization above.
On this site, related combinatorics appears in
\hyperref[temperleyLiebCrystal]{Temperley--Lieb crystals}.

\subsection[cliffordAlgebra]{Clifford algebras}

Let $V$ be a finite-dimensional vector space over a field $k$ of
characteristic different from $2$, and let $Q:V\to k$ be a quadratic form.
The \defin{Clifford algebra} is
\[
 \operatorname{Cl}(V,Q)\coloneqq
 T(V)/\langle u\otimes u-Q(u)1:u\in V\rangle,
\]
where $T(V)=\bigoplus_{r\geq0}V^{\otimes r}$ is the tensor algebra.
Writing $B(u,w)=\frac12(Q(u+w)-Q(u)-Q(w))$, the relations become
$uw+wu=2B(u,w)1$. We use the convention $u^2=Q(u)1$;
some authors use its negative. When $Q=0$, this is the exterior algebra.
In all cases, $\dim_k\operatorname{Cl}(V,Q)=2^{\dim_k V}$.

For example, over $\setC$ with $Q(x)=x_1^2+\dotsb+x_n^2$, we obtain
the algebra $\operatorname{Cl}_n$ with generators $c_1,\dotsc,c_n$ satisfying
\[
 c_i^2=1,\qquad c_ic_j=-c_jc_i\quad(i\ne j).
\]
It has basis $c_{i_1}\dotsm c_{i_r}$ with $i_1<\dotsb<i_r$, including $1$.
Declaring the generators odd makes it a $\setZ/2\setZ$-graded algebra.
The symmetric group acts by permuting the generators; the semidirect product
$\operatorname{Cl}_n\rtimes\setC[\symS_n]$ is the
\defin{Sergeev algebra}, also called the Hecke--Clifford algebra.
It connects spin representations of symmetric groups with
\hyperref[schurQ]{Schur's $Q$ functions}; see \cite[Secs.~2--3]{WanWang2012Spin}.

\subsection[doubleAffineHeckeAlgebra]{Double affine Hecke algebras}

The \defin{double affine Hecke algebra} (DAHA), introduced by
\name{Ivan Cherednik}, extends an affine Hecke algebra by a second family
of commuting invertible elements; see \cite{Cherednik1995} and
\href{https://math.mit.edu/~etingof/rivercher.pdf}{P.~Etingof's overview}.
In the type $GL_n$ version, it contains the finite Hecke generators
$T_1,\dotsc,T_{n-1}$ and two Laurent polynomial subalgebras generated by
$X_1^{\pm1},\dotsc,X_n^{\pm1}$ and $Y_1^{\pm1},\dotsc,Y_n^{\pm1}$.
The $X_i$ commute with one another, as do the $Y_i$, but the two families
do not commute with each other. Their cross-relations depend on parameters
$q,t$.

In the polynomial representation on Laurent polynomials, $X_i$ acts by
multiplication by $x_i$, while the $Y_i$ act by commuting
$q$-difference-reflection operators. For generic parameters,
\hyperref[macdonaldE]{nonsymmetric Macdonald polynomials} are their joint
eigenfunctions. Symmetric polynomials in the $Y_i$ preserve symmetric
Laurent polynomials and have the
\hyperref[macdonaldP]{symmetric Macdonald polynomials} as joint eigenfunctions.

\subsection[klyachkoAlgebra]{The $q$-Klyachko algebra}

The $q$-Klyachko algebra was introduced by \name{Philippe Nadeau} and
\name{Vasu Tewari} \cite{NadeauTewari2025QKlyachko}.  It is a
$q$-deformation of the $\symS_n$-invariant part of the rational cohomology
ring of the permutahedral variety, designed to interpolate between Klyachko's
formula and \name{Ian Macdonald}'s reduced-word formula.

It is given by 
\[
   \setQ[\dotsc,u_{-1},u_0,1,u_2,\dotsc] / \langle (q+1)u_i^2 = q u_i u_{i-1} + u_i u_{i+1} \rangle.
\]

\subsection[clusterAlgebra]{Cluster algebras}

\defin{Cluster algebras}\label{clusterAlgebraDefinition} are commutative
algebras generated from overlapping clusters by mutation, introduced by
\name{Sergey Fomin} and \name{Andrei Zelevinsky} \cite{FominZelevinsky2002}.

\subsection[hopfAlgebra]{Hopf algebras}

\defin{Hopf algebras}\label{hopfAlgebraDefinition} are algebras equipped with
compatible coproduct, counit, and antipode maps; for a combinatorial survey,
see \name{Darij Grinberg} and \name{Victor Reiner} \cite{GrinbergReiner2014Hopf}.

\name{Farid Aliniaeifard} and \name{Stephanie van Willigenburg} introduce
\defin{rooted Hopf algebras}, graded connected Hopf algebras equipped with a
group-like element in their completion, and identify their category as dual
to the category of combinatorial Hopf algebras
\cite{AliniaeifardWilligenburg2026x}.  Roots extend Bernstein and raising
operators, producing analogues of Schur, Hall--Littlewood, and Macdonald
functions; the construction is developed explicitly for symmetric functions
in noncommuting variables.
